Esta secci\'on presenta las salidas visuales generadas por el propio proyecto en Python. Las figuras se insertan directamente desde \texttt{outputs/figuras/}, de modo que el documento queda alineado con la corrida actual.

\section*{Resumen ejecutivo del portafolio}

\begin{figure}[H]
\centering
\includegraphics[width=0.95\textwidth]{01_tabla_resumen_portafolio.png}
\caption{Resumen del portafolio para la corrida actual.}
\end{figure}

La tabla anterior resume las magnitudes centrales del ejercicio. El valor esperado a 1 a\~no es mayor que el valor actual, por lo que la p\'erdida esperada aparece con signo negativo. Eso refleja que, en promedio, el efecto de carry y cupones domina a la p\'erdida esperada por migraci\'on crediticia.

\begin{figure}[H]
\centering
\includegraphics[width=\textwidth]{02_tabla_resumen_bonos.png}
\caption{Resumen comparativo de los tres bonos.}
\end{figure}

La tabla por bono deja ver una diferencia importante en perfil de riesgo: el bono A- es mucho m\'as estable, mientras que los bonos BBB- y B- concentran la mayor parte del riesgo extremo del portafolio.

\section*{Distribuciones marginales por bono}

\begin{figure}[H]
\centering
\includegraphics[width=\textwidth]{03_grafica_bono_a_menos.png}
\caption{Distribuci\'on de probabilidad y valor al horizonte para el bono A-.}
\end{figure}

En el bono A-, la masa de probabilidad permanece concentrada en ratings cercanos a su estado inicial. La dispersi\'on del valor al horizonte es baja y el default tiene muy baja probabilidad, por eso su volatilidad es la menor del grupo.

\begin{figure}[H]
\centering
\includegraphics[width=\textwidth]{04_grafica_bono_bbb_menos.png}
\caption{Distribuci\'on de probabilidad y valor al horizonte para el bono BBB-.}
\end{figure}

En el bono BBB- aparece una combinaci\'on m\'as delicada: el valor esperado sigue siendo alto, pero la cola de p\'erdidas crece de forma visible. Su VaR individual es elevado porque el default produce una ca\'ida marcada respecto al precio actual.

\begin{figure}[H]
\centering
\includegraphics[width=\textwidth]{05_grafica_bono_b_menos.png}
\caption{Distribuci\'on de probabilidad y valor al horizonte para el bono B-.}
\end{figure}

El bono B- es el activo m\'as riesgoso del portafolio. La probabilidad de deterioro severo y default es sustancialmente mayor, y eso se traduce en la mayor volatilidad individual del conjunto.

\section*{Distribuci\'on del portafolio}

\begin{figure}[H]
\centering
\includegraphics[width=\textwidth]{06_grafica_portafolio_distribucion.png}
\caption{PMF y CDF del portafolio: convoluci\'on vs brute force.}
\end{figure}

La superposici\'on entre ambas curvas muestra que la convoluci\'on reproduce exactamente la distribuci\'on obtenida por enumeraci\'on completa a nivel centavos. Adem\'as, el cuantil inferior del 0.1\% se ubica alrededor de \LowerTail, lo que equivale a una p\'erdida de aproximadamente \PortfolioVaR\ respecto al valor actual.

\begin{figure}[H]
\centering
\includegraphics[width=\textwidth]{07_grafica_pasos_convolucion.png}
\caption{Pasos intermedios de la convoluci\'on.}
\end{figure}

Esta figura es importante porque muestra que la distribuci\'on del portafolio se construye de forma incremental: primero se toma el soporte del primer bono, luego se suma el segundo y finalmente el tercero. El soporte crece, pero sin necesidad de mantener una tabla manual de los \(5{,}832\) escenarios.

\section*{Tablas visuales para el reporte}

\begin{figure}[H]
\centering
\includegraphics[width=\textwidth]{08_tabla_cola_portafolio.png}
\caption{Cola inferior del portafolio ordenada por valor.}
\end{figure}

La cola inferior concentra los peores valores del portafolio. Es la parte relevante para VaR y ES, porque all\'i viven los escenarios de estr\'es m\'aximo.

\begin{figure}[H]
\centering
\includegraphics[width=\textwidth]{09_tabla_peores_escenarios_bruteforce.png}
\caption{Peores escenarios obtenidos por brute force.}
\end{figure}

Los peores escenarios muestran que el riesgo extremo no proviene del bono A-, sino de eventos severos en los bonos BBB- y B-. El escenario umbral reportado por el modelo es \texttt{\ThresholdScenario}, que combina estabilidad relativa del bono A- con deterioro extremo en los bonos m\'as riesgosos.

\begin{figure}[H]
\centering
\includegraphics[width=\textwidth]{10_tabla_comparacion_convolucion_vs_bruteforce.png}
\caption{Comparaci\'on fila a fila entre convoluci\'on y brute force en la cola.}
\end{figure}

La comparaci\'on visual de la cola sirve como evidencia metodol\'ogica: la diferencia entre ambas PMF es nula o despreciable al nivel de redondeo usado por el proyecto.
